\sscp{Verbal and nominal morphology}{13-03-01}
\ssscp{Verbal prefixes}{13-03-01-01}
Retentions of the verbal prefixes *ma- ‘stative’,
morphological causative *pa- ‘causative’,
and PCEMP **pa(Ra)- ‘reciprocal’\footnote{
		PCEMP **paRa- ‘reciprocal
		prefix’ is from \citet{BlustTrussel2020},
		although forms in the region that support anything beyond *pa-
		‘reciprocal’ are extremely scarce.
		\citet{BlustTrussel2020} cite only \ili{Rembong} \it{para} ‘mutual,
		reciprocal’ from the target region,
		plus Tolai \it{wara-} (an Oceanic language) in support of their fuller PCEMP reconstruction.
		Varieties of \ili{Geser} on the Seram mainland have the form \it{far-}
		‘reciprocal’ \citep{Loski-Grimes1990}.
		Precursors and related forms are found widely west
		and north of the Blust line in western Indonesia and the Philippines.
		For example, one function of \ili{Malay} \it{bər-} is ‘reciprocal’,
		as in \it{bər-hadap-an} ‘face each other’,
		\it{bər-aku{\tl}aku-an} ‘use the form “\it{aku}” when talking
		to each other’, \it{bər-apit} ‘be wedged in against each other’,
		\it{bər-baik-(an)} ‘be in a good relationship with each other’,
		\it{bər-balas-an} ‘exchange (responses, accusations,
		shots) with each other’,
		\it{bər-bariŋ-an} ‘lie next to each other’, etc.
		\citet{BlustTrussel2020} have a note under PWMP *paR-1 saying
		“\citet{Wolff1981} considers the form
		and meaning of *paR- within the context of a larger comparison
		of the morphology of \ili{Tagalog} and Bahasa Indonesia.
		He concludes (p. 87, fn. 14): ‘The meanings of \it{pag-}
		and \it{per-} are elusive and range widely.
		We doubt that they can be profitably compared.’ It is quite possible
		that *paR- had multiple functions in PWMP,
		and that the meaning reconstructed here is only one of several
		glosses that may eventually be associated with this etymon.”}
are extremely widespread in many or all subgroups in the region.
They are productive in some languages and
nproductive in others occurring only as frozen forms.
The causative *pa- and reciprocal *pa(Ra)- often surface as homophones
in the daughter languages in the region
(\citealp[113--115]{Grimes1991c}, \citealp{Grimes2008,GrimesRanohAplugi2008}).

Comparative data for retentions of these three verbal prefixes
are shown in \trf{13-12}.
Reflexes in some subgroups is relatively easy to find,
and quite sparse in others due to the scarcity of grammatical descriptions.
Reciprocals in many languages are indicated by manipulating \tsc{3pl}
pronouns in various ways (e.g.
\it{ra-si}-verb, \it{ra}-verb \it{sir}), or by a separate word,
rather than by a verbal prefix.
Many languages have both morphological causatives and periphrastic causatives.
In \trf{13-09}, *(p)a- indicates some investigators describing
languages within that subgroup have identified \it{pa-} as a
frozen or archaic form in one or more items,
but the normal reflexes are \it{fa-,
ha-, a-}, according to the regular sound correspondences for
*p for that language
and subgroup (see \trf{13-02}).
Words with no asterisk are forms in individual languages.

\begin{table}
	\centering\tcp{Some verbal prefixes in Wallacean subgroups (preliminary)}{13-12}
	\st{6pt}\begin{tabular}{llll}\lsptoprule
Subgroup						&	Stative	&	Causative	&	Reciprocal\\	\midrule
\tsc{\ili{Bima}-Lembata}	&	*ma-	&	*pa-	&	*paRa-\\
\tsc{Timor-Babar}		&	*ma-	&	*(p)a-	&	{\hr}\it{apa-, hak-, wa-}\\
\tsc{\ili{Central} Timor}	&	*ma-	&	*a-	&	*a- (?)\\
\tsc{Aru}						&	{\hr}\it{ma-}	&	{\hr}\it{ka-}	&	{\hr}\it{r-}\\
STB									&	*ma-	&	*pa-	&	*paR-\\
\ili{Banda}								&	\it{-m(a)-, mo-}	&	{\hr}\it{mbo-, (fo-)}	&	\\
\tsc{\ili{Ambon}-Seram}		&	*ma-	&	*pa-	&	\it{\cg{(baku)}}\\
\tsc{Sula-Buru}			&	*ma-	&	*pa-	&	*pa-\\
		\lspbottomrule
	\end{tabular}
\end{table}

Reflexes of stative *ma- tend to be associated with non-active
verbs, some experiencer verbs, and concepts,
qualities or properties that are considered “adjectives” in many Indo-European languages.
The following examples are from the \tsc{\ili{Bima}-Lembata} language \ili{Dhao}.
Example \qf{13-19} illustrates a \it{ma-} ‘stative’ prefix in
\ili{Dhao}, marking the subject as Undergoer rather than Actor,
and in a BE relationship rather than a DO relationship with the
quality, state, or process of the non-active verb.
\ili{As} explained in \srf{03-03-01},
penultimate schwa in \ili{Dhao} is syllabic,
stress bearing, and causes phonetic lengthening of the following consonant.
We also described how *uCa > \it{əCu}
and *iCa > \it{əCi} in \ili{Havu} and \ili{Dhao}.
(Examples are from \citealp{GrimesRanohAplugi2008}.)

\noindent{\wg\st{6pt}\begin{tabularx}
{\linewidth}{>{\hspace{-\tabcolsep}}l<{\hspace{-\tabcolsep}}
>{\hspace{-\tabcolsep}\enspace}ll
>{\raggedright\arraybackslash}X}
\lsptoprule
%\tbx{13-19}	&	\ili{Dhao}	&	\ili{Dhao} gloss			&		\\	\midrule
%\endfirsthead
\tbx{13-19}	&	\ili{Dhao}	&	\ili{Dhao} gloss	&		\\	\midrule
%\endhead
	&	ma-ɖ͡ʐaʔu	&	be afraid, fearful&	PMP *ma-takut `fearful, afraid'	\\
	&	ma-ɖ͡ʐasa	&	be ripe (fruit)		&	PMP *ma-tasak `ripe, cooked'	\\
	&	ma-ɖ͡ʐera	&	be long						&	\hp{\,}	\\
	&	ma-ɗəka		&	be sharp					&	\hp{\,}	\\
	&	ma-ʄəni		&	be diligent				&	\hp{\,}	\\
	&	ma-kae		&	be ashamed				&	also:	\it{pa-ma-kae} `cause to be ashamed, shame s.o.'	\\
	&	ma-laa		&	be amazed					&	also:	\it{pa-ma-laa} `cause to be amazed, amaze s.o.'	\\
	&	ma-lai		&	be quick					&	\hp{\,}	\\
	&	ma-ŋaŋa		&	be hungry					&	PMP *ŋaŋa open the mouth wide, gape; gaping	\\
	&	ma-ɲilu		&	be sour						&	PMP *ñilu painful sensation in teeth (from eating s.t. sour)	\\
	&	ma-rəŋi		&	be calm (of sea)	&	PMP *linaŋ calm, tranquil, of the surface of water	\\
	&	ma-rəma		&	be deep						&	\tsc{PAn} *ma-daləm deep	\\
\lspbottomrule
\end{tabularx}\mbox{}\\}

Example \qf{13-20} shows the causative prefix \it{pa-} in \ili{Dhao}
and example \qf{13-21} shows the homophonous reciprocal prefix
\it{pa-} in \ili{Dhao}.\footnote{
		\citet[9]{Balukh2020} seems unaware
		of this very common pattern of homophony in the region
		and treats the causative
		and reciprocal functions as a single prefix,
		“\ili{Dhao} has only one derivational affix; that is the prefix is \it{pa-}.
		It is used to derive verbs from nouns
		and adjectives, as well as change the valence of verbs.
		Semantically, the prefix \it{pa-} expresses causative,
		reciprocal, intensity, and other meanings.”}

\noindent{\wg\st{6pt}\begin{tabularx}
{\linewidth}{>{\hspace{-\tabcolsep}}l<{\hspace{-\tabcolsep}}
>{\hspace{-\tabcolsep}\enspace}llll
>{\raggedright\arraybackslash}X}
\lsptoprule
%\tbx{13-20}	&	\ili{Dhao}	&	gloss	&	\ili{Dhao}	&	gloss	\\	\midrule
%\endfirsthead
\tbx{13-20}	&	\ili{Dhao}	&	gloss	&	\ili{Dhao}	&	gloss	\\	\midrule
%\endhead
	&	ɓaɓa	&	short				&	pa-ɓaɓe		&	\mc{2}{l}{shorten s.t.}					\\
	&	beʔa	&	good				&	pa-beʔa		&	\mc{2}{l}{fix s.t.}							\\
	&	b͡βəlu	&	forget			&	pa-b͡βəlu	&	\mc{2}{l}{forget s.t. (actively, deliberately)}	\\
	&	ʧəna	&	sink, drown	&	pa-ʧəne		&	\mc{2}{l}{drown s.o./s.t.} 										\\
	&	deɖ͡ʐa	&	above				&	pa-deɖ͡ʐa	&	\mc{2}{l}{elevate s.t., lift up, put up}			\\
	&	haha	&	low, below	&	pa-haha		&	lower s.t., put down 						&	(< PMP *babaq low, below)	\\
	&	huni	&	hide				&	pa-huni		&	hide s.o./s.t. 									&	(< PMP *buni hide, conceal)	\\
	&	ʄunu	&	sleep				&	pa-ʄunu		&	\mc{2}{l}{put to sleep, cause to lie down}						\\
	&	əʧi		&	one					&	pa-ʔəʧi		&	join, unite 										&	(< PMP *isa one)		\\
	&	ŋara	&	name				&	pa-ŋare		&	name s.o./s.t. 									&	(< PMP *ŋajan name)	\\
\lspbottomrule
\end{tabularx}\mbox{}\\}

\noindent{\wg\st{6pt}\begin{tabularx}
{\linewidth}{>{\hspace{-\tabcolsep}}l<{\hspace{-\tabcolsep}}
>{\hspace{-\tabcolsep}\enspace}lll
>{\raggedright\arraybackslash}X}
\lsptoprule
%\tbx{13-21}	&	\ili{Dhao}	&	gloss	&	\ili{Dhao}	&	gloss	\\	\midrule
%\endfirsthead
\tbx{13-21}	&	\ili{Dhao}	&	gloss	&	\ili{Dhao}	&	gloss	\\	\midrule
%\endhead
	&	bala	&	respond	&	pa-bala	&	discuss, argue with each other	\\
	&	ɗəi	&	like, want	&	pa-ɗəi	&	mutual like, like each other	\\
	&	game	&	hit, beat up	&	pa-gama	&	fight each other	\\
	&	karəi	&	ask	&	pa-karəi	&	ask each other, question each other	\\
	&	kasiro	&	shoot weapon	&	pa-kasiro	&	shoot each other	\\
	&	laɓa	&	oppose	&	pa-laɓa	&	oppose each other	\\
	&	peka	&	speak	&	pa-peka	&	talk to each other, tell each other	\\
	&	taŋara	&	face s.o.	&	pa-taŋara	&	face each other	\\
	&	too	&	border	&	pa-too	&	share a border with each other	\\
	&		&		&	pa-ʧəri	&	separate from each other, divorce (cf. \ili{Malay} \it{cərai})	\\
	&		&		&	pa-katoo	&	curse each other	\\
\lspbottomrule
\end{tabularx}\mbox{}\\}

\ili{As} retentions rather than innovations,
these historical prefixes are not useful for subgrouping purposes.
Although as noted in \srf{04-02},
at a lower level \tsc{Rote-Meto} languages appear to have locally
innovated a prefix \it{ma-} ‘reciprocal’,
which is consequently homophonous with \it{ma-} ‘stative’,
and \it{ma-} ‘have (characteristic, quality)’.
The \tsc{Rote-Meto} prefix \it{ma-} ‘reciprocal’ does not appear
to be found in other nodes of the \tsc{Timor-Babar} subgroup.

\ssscp{Applicatives}{13-03-01-02}
There is also a fairly widespread verbal applicative that often
takes the form \it{-aka,
-ak, -ik, -k, -ke, -ko, -ʔ, -ŋ}.
Some of the functions documented for these include: deriving
verbs from nouns, transitivising intransitive verbs, forming causatives,
shifting aspect from activity to accomplishment interpretation,
and various kinds of repackaging the valence
and role structure of the verb such as dative raising
and argument incorporation, or simply flagging the end of the
clause nucleus by marking the last element considered to be verbal
(see \citealp{Jonker1906,Grimes1991b}, \citealp[108-112]{Grimes1991c}, \citealp{Tamelan2021}).
These are related historically to the form reconstructed for
P(W)MP *akən ‘preposition, on, upon; applicative marker’ that is a suffix in many languages
west of the Blust line,
such as \ili{Malay} \it{-kan},
and \ili{Tukang Besi} \it{-ako} \citep{Donohue1999a},
and share a similar range of functions as those described for
western Austronesian languages \citep{TruongMcDonnell2021}.
Due to the severe under-documentation in so many languages in
the region, we simply note its presence
and its similarities in form
and range of functions.
There is much room for further documentation and comparative study.
\ili{As} a retention rather than an innovation,
this historical applicative suffix/enclitic is not useful for subgrouping purposes.

\ssscp{Infix *\<in\>}{13-03-01-03}
There is evidence of scattered reflexes of a historical infix
*\<in\> ‘marker of deverbal nouns’ in \tsc{Sula-Buru},
\tsc{Seram-Tanimbar-Bomberai}, and \tsc{Timor-Babar} (especially \tsc{Southwest Maluku}.\footnote{
		\citet[575]{SchapperZobel2024a} link these reflexes to PMP *\<in\> “inflectional
		past/perfective TAM marker”, but the pathway from a verbal ‘perfective marker’ to marking
		‘deverbal nouns’ is not immediately transparent.
		\citet{BlustTrussel2020} list nominalizing functions for the
		\tsc{PAn} infix *\<in\> ‘perfective marker’ for languages west of the Blust line
		\ili{Wolio} and \ili{Chamorro}, as well as for Oceanic languages Tolai, Hoava and Roviana.
		Their note also acknowledges that \citet{StarostaPawleyReid1982} saw this
		and similar affixes as “derivational devices for creating deverbal nouns”.
		\citet{BlustTrussel2020} see the possibility that \tsc{PAn} *\<in\> could be both
		verbal and nominal, rather than either verbal or nominal.
		\citet[163]{StarostaPawleyReid1982} see the nominal function
		at the level of \tsc{PAn} as prior to
		and more basic than the verbal perfective function saying,
		“The affix *ni-/-in- functioned in Proto-Austronesian to derive
		nouns from verbs and other nouns, although it may have also begun to have the
		function of marking perfective aspect in verbs,
		a function which is now its primary one in Philippine languages.”
		Seeing the forms in Wallacea that derive deverbal nouns as a
		retention of \tsc{PAn} *\<in\> deriving deverbal nouns is straightforward,
		whereas deriving the nominalizing functions in Wallacea from
		a PMP ‘perfective marker’ that developed later in Philippine languages is convoluted.
		There is certainly room for a deeper delve into this infix (or
		these infixes) comparatively and historically.
		But as a retention,
		these traces in Wallacea are not relevant for subgrouping.}
Synchronically reflexes of *\<in\> appear as prefixes,
infixes, or trigger metathesis in Wallacea.
Some examples adapted from \citet[575f]{SchapperZobel2024a} are
shown in example \qf{in1}.
A similar nominalising infix is also reported for the \tsc{Timor-Babar}
languages \ili{Leti} and \ili{Luang}.
\\

\noindent{\wg\st{6pt}\begin{tabularx}
{\linewidth}{>{\hspace{-\tabcolsep}}l<{\hspace{-\tabcolsep}}
>{\hspace{-\tabcolsep}\enspace}llll
>{\raggedright\arraybackslash}X}
\lsptoprule
%\tbx{(AA)}	&	Subgroup	&	Language	&	Morpheme	&	Verb	&		&	Deverbal noun	\\	\midrule
%\endfirsthead
\tbx{in1}		&	Subgroup	&	Language&	Morpheme			&								& Gloss\\	\midrule
%\endhead
						&	TB				&	\ili{Roma}		&	\<ny\>/\<ɲ\>	&	sinin 				&	sing		\\	\hhline{~}
						&						&					&								&	s\<ny\>inin 	&	song		\\
						&	TB				&	\ili{Wetan}		&	\<ini\>				&	wali					&	answer (v.)	\\	\hhline{~}
						&						&					&								&	w\<ini\>ali		&	answer (n), response	\\
						&	STB				&	\ili{Yamdena}	&	\<ny\>/\<ɲ\>	&	tasiŋ					&	cry (v.)		\\	\hhline{~}
						&						&					&								&	t\<ny\>asiŋ		&	(act of) crying	\\
						&	STB				&	\ili{Fordata}	&	\<n\>					&	vuat					&	load (v.)		\\	\hhline{~}
						&						&					&								&	v\<n\>uat			&	cargo	\\
						&						&					&								&	tabar					&	stomp (v.)	\\	\hhline{~}
						&						&					&								&	t\<n\>abar		&	stomping dance	\\
\lspbottomrule
\end{tabularx}\mbox{}\\}

\citet[65f, 145]{Grimes1991c}
describes a similar ‘participial’ prefix \it{en-} for \tsc{Sula-Buru}
language \ili{Buru}, which derives deverbal nouns (including abstract nouns) from
active verb roots only,\footnote{
		The ability (or not) to take \it{en-} is one of the morphosyntactic
		tests for distinguishing active from non-active verb roots in \ili{Buru}.}
and whose allophones are dependent on the phonological shape
of the onset of the verb root.
Before most C-initial roots,
\it{en-} is a simple prefix sometimes assimilating to point of articulation,
as in example \qf{in2}.
Prefixes added to V-initial roots in \ili{Buru} trigger palatalization,
so \it{en-} → [eɲ-] as in example \qf{in3}.
When connected to roots starting with /f k s t/,
the /en-/ triggers metathesis of the adjacent consonants as in example \qf{in4}.\footnote{
		\ili{Buru} [ʔ] is an archiphone,
		allophonic of voiceless obstruents /p t k/ in a restricted environment and
		/e/ is the default vowel for all \ili{Buru} prefixes,
		and not a historial retention.}

\noindent{\wg\st{6pt}\begin{tabularx}
{\linewidth}{>{\hspace{-\tabcolsep}}l<{\hspace{-\tabcolsep}}
>{\hspace{-\tabcolsep}\enspace}llll
>{\raggedright\arraybackslash}X}
\lsptoprule
	%&	stem	&	gloss	&		&	derivation	&	gloss	&		&	derivation	&	gloss	\\	\midrule
%\endfirsthead	
\tbx{in2} 	&	stem	&	gloss		&				&	derivation	&	gloss	\\	\midrule
%\endhead
	&	dapi		&	smoke meat			&	{\ra}	&	/en-dapi-t/	&	\tsc{ptcp}-smoke meat-\tsc{nmlz}	\\	\hhline{~}
	&					&									&				&	endapit			&	smoked (meat)	\\
	&	heka		&	flee						&	{\ra}	&	/en-heka-t/	&	\tsc{ptcp}-flee-\tsc{nmlz}				\\	\hhline{~}
	&					&									&				&	enhekat			&	escape (n.)	\\
	&	hesa-k	&	pull s.t., drag	&	{\ra}	&	/en-hesa-t/	&	\tsc{ptcp}-pull-\tsc{nmlz}				\\	\hhline{~}
	&					&									&				&	enhesat			&	influence (n.)	\\
	&	laha		&	request (v)			&	{\ra}	&	/en-laha-t/	&	\tsc{ptcp}-request-\tsc{nmlz}			\\	\hhline{~}
	&					&									&				&	enlahat			&	request (n.)	\\
	&	mata		&	die							&	{\ra}	&	/en-mata/		&	\tsc{ptcp}-die										\\	\hhline{~}
	&					&									&				&	enmata			&	death	\\
	&	reha		&	tempt, goad			&	{\ra}	&	/en-reha-n/	&	\tsc{ptcp}-tempt-\tsc{gen}				\\	\hhline{~}
	&					&									&				&	enrehan			&	temptation	\\
	&	wada		&	shoulder carry	&	{\ra}	&	/en-wada-t/	&	\tsc{ptcp}-carry-\tsc{nmlz}				\\	\hhline{~}
	&					&									&				&	enwadat			&	burden	\\
	&	eʔbana	&	give birth			&	{\ra}	&	/en-bana-t/	&	\tsc{ptcp}-birth-\tsc{nmlz}				\\	\hhline{~}
	&					&									&				&	embanat			&	birth	\\
	&	puna		&	make, do				&	{\ra}	&	/en-puna-n/	&	\tsc{ptcp}-do-\tsc{gen}						\\	\hhline{~}
	&					&									&				&	empunan			&	behaviour	\\
	&	gosa		&	do/be good			&	{\ra}	&	/en-gosa-n/	&	\tsc{ptcp}-good-\tsc{gen}					\\	\hhline{~}
	&					&									&				&	eŋgosan			&	goodness, kindness	\\
\lspbottomrule
\end{tabularx}\mbox{}\\}

\noindent{\wg\st{2pt}\begin{tabularx}
{\linewidth}{>{\hspace{-\tabcolsep}}l<{\hspace{-\tabcolsep}}
>{\hspace{-\tabcolsep}\enspace}llll
>{\raggedright\arraybackslash}X}
\lsptoprule
\tbx{in3}	&	Underlying	&	Gloss													&				&	Realisation	&	Gloss	\\	\midrule
					&	/en-oli-t/	&	\tsc{ptcp}-return-\tsc{nmlz}	&	{\ra}	&	eɲolit	&	return (trip, fare)	\\
					&	/en-opi-t/	&	\tsc{ptcp}-blow-\tsc{nmlz}		&	{\ra}	&	eɲopit	&	thunderstorm, windstorm	\\
					&	/en-ino-t/	&	\tsc{ptcp}-drink-\tsc{nmlz}		&	{\ra}	&	eɲinut	&	drinking (water)	\\
					&	/en-ego-n/	&	\tsc{ptcp}-take-\tsc{gen}			&	{\ra}	&	eɲegun	&	1) taken (goods); 2) harvest	\\
\lspbottomrule
\end{tabularx}\mbox{}\\}

\noindent{\wg\st{4pt}\begin{tabularx}
{\linewidth}{>{\hspace{-\tabcolsep}}l<{\hspace{-\tabcolsep}}
>{\hspace{-\tabcolsep}\enspace}llll
>{\raggedright\arraybackslash}X}
\lsptoprule
%(DD)	&	Underlying	&	Gloss	&		&	Realisation	&	Gloss	\\	\midrule
%\endfirsthead
\tbx{in4}	&	Underlying	&	Gloss	&		&	Realisation	&	Gloss	\\	\midrule
%\endhead
	&	/en-foi-t/	&	\tsc{ptcp}-bathe-\tsc{nmlz}	&	{\ra}	&	efnoit	&	basin	\\
	&	/en-fasa-t/	&	\tsc{ptcp}-cut-\tsc{nmlz}	&	{\ra}	&	efnasat	&	decision	\\
	&	/en-kadu-t/	&	\tsc{ptcp}-come-\tsc{nmlz}	&	{\ra}	&	eʔnadut	&	arrival	\\
	&	/en-kodo-t/	&	\tsc{ptcp}-grind-\tsc{nmlz}	&	{\ra}	&	eʔnodot	&	grinding (stone)	\\
	&	/en-suba-n/	&	\tsc{ptcp}-cross.thresh-\tsc{gen}	&	{\ra}	&	esnuban	&	its threshold	\\
	&	/en-seba-t/	&	\tsc{ptcp}-worship-\tsc{nmlz}	&	{\ra}	&	esnebat	&	worship (house)	\\
	&	/en-toho-n/	&	\tsc{ptcp}-descend-\tsc{gen}	&	{\ra}	&	eʔnohon	&	descent (social)	\\
	&	/en-tati-n/	&	\tsc{ptcp}-set.down-\tsc{gen}	&	{\ra}	&	eʔnatin	&	1) tribute, tax; 2) genealogy	\\
\lspbottomrule
\end{tabularx}\mbox{}\\}

There may also be traces of an infix *\<in\> in the
\tsc{Flores-Lembata} node of \tsc{\ili{Bima}-Lemabata}.
There are several words which
occur with medial \it{n}.
Examples from \citet{Fricke2019} include:
*salaq > \ili{Kalikasa} \it{s\tbr{n}alak} `wrong',
*ma-həyaq > \ili{Alorese} \it{m\tbr{n}iaŋ} `shy, ashamed',
and \ili{Sanskrit} \it{sama} > \ili{Kalikasa} \it{s\tbr{n}amaŋ} `same'.

These languages synchronically all have either infixing or trigger
some metathesis, and collectively suggest a historical infix *\<in\> ‘deriving deverbal nouns’.
This proto-form and its reflexes warrants further study.

\ssscp{Nominalisers}{13-03-01-04}
There are also some fairly widespread nominalisers that often
take the forms \it{-t, -te, -ti, -s, -k, -ke, -ʔ} in the region.
Common functions in geographically distant languages (such as
in \tsc{Sula-Buru}, \ili{Taliabu}, \tsc{\ili{Ambon}-Seram}, and  \tsc{Timor-Babar} (both in \tsc{Rote-Meto}
and \tsc{Southwest Tanimbar}) include: making nouns out of verbs
or precategorials, and marking a verb as being used attributively within a nominal
constituent in the clause (an NP)
(\citealp{Jonker1906}, \citealp[30--33]{Collins1983}, \citealp[101--104, 143--145, 180--182]{Grimes1991c}, \citealp{Tamelan2021}).
These can also attach to body parts
and “inalienable nouns”, indicating they are independent,
unpossessed, amputated, disembodied, generic,
unassociated with or independent of the “whole” in a physical
or conceptual part-whole relationship.
This latter function is something like a zero person (unassociated, ≈obviative)
genitive (\citealp[442]{Edwards2020}, \citealp{Tamelan2021}).
The \it{-f} suffix in \tsc{Meto} languages often has this latter
function, as does \it{-r} in \ili{Kemak}.
Some also have a discourse function tracking parts of wholes
whose identity has already been established in the discourse,
or whose “whole” is assumed to be known from the larger context
or from general knowledge,
such that the genitive part-whole relationship no longer needs
to be indicated (\citealp[288f]{Grimes1991c}, \citealp[617f]{Grimes2024}).

This nominaliser shows possibilities for subgrouping as a morphological
innovation, particularly since similar forms,
functions and distributions of \it{-t} are found in \tsc{Sula-Buru},
\tsc{\ili{Ambon}-Seram} and also in \tsc{Rote-Meto} languages,
but apparently not in most of the wider \tsc{Timor-Babar} subgroup
to which the \tsc{Rote-Meto} languages belong.
Given this odd distribution,
parallel development cannot be ruled out.
Given the severe under-documentation of the languages in the
region, it is as yet unclear comparatively how widespread possible cognates
might be, including outside the region.
It is also unclear if this nominaliser developed from an earlier
article, case marker, deictic, or some other particle.
This topic is clearly worthy of further study,
but the work needed to do an adequate job is beyond the scope
of this present study.

\ssscp{Relativisers}{13-03-01-05}
Various relativisers are found throughout the region.
It is possible that scattered relativisers throughout Wallacea
shaped \it{ma, mak, man, maŋ, ka} relate to a common form such as **maka,
**maŋa, or **mana ‘relativiser’, or ultimately are retentions
of something like \tsc{PAn} *manu ‘interrogative marker: which?’.
They are interspersed with non-cognate relativisers of many other forms (see \trf{13-10}).

Broadly speaking, there are three types of relativisers in the region.
The first is the common relativiser for marking relative clauses; ‘which, who’.
The second is often derived from the first for contrastive focus
constructions ‘the one who did X (in contrast to someone else;
labelled ‘emphatic’ by some analysts)’.
The third is usually restricted to (derived) nominals
and involves \it{characteristic roles} of a profession,
activity or quality associated with the identity of a person,
much like \ili{English} \it{-er} (\it{farmer,
carpenter, driver, dreamer, murderer}) -- the one who does X or
is characterised as X.
In some languages this latter function is handled grammatically
with a special relativiser,
whereas in other languages the equivalent function takes other
nominal morphology, and in still others it is handled by simple juxtaposition,
as in \it{person + work = worker}.
Relativisers flagging role are illustrated for the \tsc{\ili{Ambon}-Seram}
language South \ili{Nuaulu} \citep[61--63,
77f, 164--170]{Bolton1990} in example \qf{13-22},
and in example \qf{13-23}
for the \tsc{Sula-Buru} language of \ili{Buru} \citep[183, 429--436]{Grimes1991c},
using the habitual relativiser \it{ka} with the noun \it{geba} ‘person’.

\noindent{\wg\st{6pt}\tblr{>{\hspace{-\tabcolsep}}l<{\hspace{-\tabcolsep}}
>{\hspace{-\tabcolsep}\enspace}
llll}
\lsptoprule
%\tbx{13-22}	&	\ili{Nuaulu}	&	gloss	&	relativised	&	gloss	\\	\midrule
%\endfirsthead
\tbx{13-22}	&	\ili{Nuaulu}	&	gloss	&	relativised	&	gloss	\\	\midrule
%\endhead
	&	rahu	&	plant (v.)	&	ia \tbr{mam-}rahu-e	&	farmer	\\
	&	akarota	&	lie	&	ia \tbr{m-}akarota-ne	&	liar	\\
	&	atanunu	&	cook (v.)	&	ia \tbr{m-}atanunu-e	&	a cook (n.)	\\
	&	ehu	&	smoke (v.)	&	ia \tbr{mam-}ehu-e	&	smoker	\\
\lspbottomrule
\end{tabular}\mbox{}\\}

\noindent{\wg\st{6pt}\begin{tabularx}
{\linewidth}{>{\hspace{-\tabcolsep}}l<{\hspace{-\tabcolsep}}
>{\hspace{-\tabcolsep}\enspace}lll
>{\raggedright\arraybackslash}X}
\lsptoprule
%\tbx{13-23}	&	\ili{Buru}	&	gloss	&	relativesd	&	gloss	\\	\midrule
%\endfirsthead
\tbx{13-23}	&	\ili{Buru}	&	gloss	&	relativesd	&	gloss	\\	\midrule
%\endhead
	&	ruba	&	treat	&	geba \tbr{ka} ruba geba	&	healer, medical practitioner	\\
	&	epmata	&	kill	&	geba \tbr{ka} epmata geba	&	murderer, killer of people	\\
	&	fasa	&	cut, decide	&	geba \tbr{ka} fasa perkara	&	judge, decider of disputes [<\ili{Sanskrit}]	\\
	&	lata	&	cut	&	geba \tbr{ka} lata hawa	&	farmer, cutter of fields	\\	\hhline{~}
	&				&		&	geba \tbr{ka} lata kau	&	carpenter, cutter of wood	\\
	&	toto	&	pound	&	geba \tbr{ka} toto momol	&	blacksmith, pounder of iron	\\
	&	eʔnilik	&	exchange	&	geba \tbr{ka} eʔnilik liet	&	translator, interpreter, exchanger of language	\\
\lspbottomrule
\end{tabularx}\mbox{}\\}

For \ili{Buru},
\it{ha} is the general relativiser which identifies information
assumed to be known that enables the addressee to identify the correct referent.
Because the information in the restrictive relative clause is
assumed to be known,
the relative clause with \it{ha} is often bracketed with a definite
deictic \it{dii} ‘distal’ or \it{naa} ‘proximal’.
The head noun of the argument in the main clause has a coreferential
generic noun or pronoun within the relative clause.
This is illustrated in example \qf{13-24}.
The relative clause identifies the two brothers in an Object
NP that got burned as those that were killed,
in contrast with the two brothers that did the killing.

\begin{exe}\setcounter{xnumi}{23}
	\ex{\gll	Petu du pefa hum.tapa dii ep-sia-k tu {\brac} kaka-wai-t rua \tbf{ha} dii du ba ep-mata dii \bracr{NP}\\
						\tsc{seq} \tsc{3pl} burn house.smoke \tsc{dist} \tsc{caus}-one-\tsc{appl} with {}
						elder-younger-\tsc{nmlz} two \tsc{rel} \tsc{dist} \tsc{3pl} \tsc{dur} \tsc{caus}-die \tsc{dist} {}\\}
			\glt	‘Then they burned down the hunting lodge together with [those two brothers whom they had killed.]’ \jambox*{(\ili{Buru}: relative clause)}
						\label{ex:13-24}
\end{exe}

The \ili{Buru} relativiser can be marked according to the plurality
of the referent, as \it{ha} ‘singular referent’ or \it{ha-r} ‘plural referent’,
as in example \qf{13-25}.

\begin{exe}
	\ex{\gll	Petu geba \tbf{ha-r} dii, ba poto nunu-r ʧeŋke-r dii, {\ldots}\\
						\tsc{seq} person \tsc{rel-pl} \tsc{dist} \tsc{dur} burn \tsc{3pl.poss-pl} clove-\tsc{pl} \tsc{dist} {}\\}
			\glt	‘Then those people whose cloves got burned, {\ldots}’ \jambox*{(\ili{Buru}: relative clause)}
						\label{ex:13-25}
\end{exe}

However, aside from those relativisers that might be retentions,
the diversity of non-cognate forms within subgroups
and throughout the region illustrated in \trf{13-13},
makes it difficult to see subgrouping possibilities based on
relativisers above the local level.
(In \trf{13-13} (\it{yaŋ}) indicates a \ili{Malay} loan.)

\begin{table}[p]\centering\tcp{Diversity of forms for relativisers in Wallacea}{13-13}
%\begin{adjustbox}{totalheight=\textheight}
\begin{adjustbox}{totalheight=\textheight}
%\begin{adjustbox}{width=\textwidth}
	%\st{2pt}\begin{tabularx}{\linewidth}{llll>{\raggedright\arraybackslash}X}\lsptoprule
	\wg\st{5pt}\begin{tabular}{lllll}\lsptoprule
	Language	&	Form	&	Function	&	Subgroup	&	Comment	\\	\midrule
	\ili{Ngadha}	&	da, (yaŋ)	&	\tsc{rel}	&	\tsc{BL}	&	\citet[44, 100, 145f]{Djawanai1983}	\\	
	\ili{Keo}	&	ta	&	\tsc{rel}	&	\tsc{BL}	&	\citet[409--412]{Baird2002}	\\	
	\ili{Lio}	&	eo	&	\tsc{rel}	&	\tsc{BL}	&	\citet{Elias2018}	\\	
	\ili{Solor} \ili{Lamaholot}	&	pe	&	\tsc{rel}	&	\tsc{BL}	&	\citet[196ff]{Kroon2016}	\\	
	\ili{Kalikasa}	&	(yaŋ)	&	\tsc{rel}	&	\tsc{BL}	&	\citet[77]{Fricke2019}	\\	
	\ili{Laboya}	&	ha-, ka-	&	\tsc{rel}	&	\tsc{BL}	&	\citet{Verdizade2019a}	\\	
	\ili{Kambera}	&	ma-, pa-	&	\tsc{rel}	&	\tsc{BL}	&	\citet[316f]{Klamer1998a}	\\	
	\ili{Dhao}	&	ɖ͡ʐu	&	\tsc{rel}	&	\tsc{BL}	&	\citet{GrimesRanohAplugi2008}	\\	\hhline{~}
		&	ka ɖ͡ʐu	&	contr. focus	&		&		\\	
	\ili{Havu}	&	do	&	\tsc{rel}	&	\tsc{BL}	&	\citet{Grimes2008}	\\	
	\ili{Dela}	&	mana, fo	&	\tsc{rel}	&	\tsc{TB}	&	\citet[401ff]{Tamelan2021};\\		\hhline{~}
				&						&						&						&	only \it{fo} for non-subject arguments	\\	
%	\ili{Bilbaa}	&	lia	&	\tsc{rel}	&	\tsc{TB}	&	\citet[293]{Edwards2021}	\\	
	Amfo{\Q}an	&	leʔ	&	\tsc{rel}	&	\tsc{TB}	&	\citet{GrimesSuanaOraCulhane2021}	\\	\hhline{~}	
		&	es leʔ	&	contr. focus	&		&		\\		\hhline{~}
		&	a- {\ldots} -t	&	role, \tsc{nmlz}	&		&	not \tsc{rel}	\\	
	Roi{\Q}is \ili{Amarasi}	&	heʔ	&	\tsc{rel}	&	\tsc{TB}	&	Edwards and Grimes fieldnotes	\\		\hhline{~}
		&	es heʔ	&	contr. focus	&		&		\\		\hhline{~}
		&	ka- {\ldots} -t	&	role, \tsc{nmlz}	&		&	not \tsc{rel}	\\	
	\ili{Helong}	&	in	&	\tsc{rel}	&	\tsc{TB}	&	\citet{UBB2012a}; 	\\	\hhline{~}	
		&	man	&	role	&		&	Grimes fieldnotes	\\	
	\ili{Tetun}	&	mak, ma,	&	\tsc{rel}, role	&	\tsc{TB}	&	\citet[70-78]{vanKlinken1999}; Grimes fieldnotes;	\\		\hhline{~}
		&	nebee	&	\tsc{rel}	&		&	 (variant: \it{maka})	\\	
	\ili{Galolen}	&	be,	&	\tsc{rel}	&	\tsc{TB}	&	\citet{CristoReiGrimes2011}	\\		\hhline{~}
		&	mak	&	role	&		&		\\	
	Dadu{\Q}a	&	namee	&	\tsc{rel}	&	\tsc{TB}	&	\citet{Penn2006}	\\	
	\ili{Raklungu}	&	maʔ, manei	&	\tsc{rel}, role	&	\tsc{TB}	&	Grimes fieldnotes	\\	
	\ili{Luang}	&	maka k-	&	\tsc{rel}	&	\tsc{TB}	&	\citet[117ff]{TaberTaber2015}	\\		\hhline{~}
		&	ha	&	\tsc{rel}	&		&		\\	
	\ili{Selaru}	&	ma-	&	\tsc{rel}	&	\tsc{TB}	&	\citet{Coward2005}	\\	
	\ili{Kemak}	&	naba ta	&	\tsc{rel}	&	\tsc{CT}	&	Celcia Elliot (pers. comm.)	\\	
	S. \ili{Mambae}	&	bae ke,	&	\tsc{rel}	&	\tsc{CT}	&	\citet{GrimesMarcalFereira2014}	\\		\hhline{~}
		&	fe, ke	&	\tsc{rel}	&		&		\\	
	\ili{Sanana}	&	(yaŋ)	&	\tsc{rel}	&	\tsc{SB}	&	\citet[330]{Bloyd2020}	\\	
	\ili{Buru}	&	ha	&	\tsc{rel}	&	\tsc{SB}	&	\citet[183, 429--436]{Grimes1991c}	\\		\hhline{~}
		&	ka	&	role	&		&		\\	
	\ili{Larike}	&	(yaŋ)	&	\tsc{rel}	&	\tsc{AS}	&	\citet{LaidigLaidig1991}	\\	
	\ili{Asilulu}	&	(yaŋ)	&	\tsc{rel}	&	\tsc{AS}	&	\citet[63]{Collins2003a}	\\		\hhline{~}
		&	maka-	&	role	&		&		\\	
	\ili{Sou Amana Teru}	&	(yaŋ)	&	\tsc{rel}	&	\tsc{AS}	&	\citet{Musgrave2010}	\\	
	\ili{Yalahatan}	&	re	&	\tsc{rel}	&	\tsc{AS}	&	Beckley n.d.	\\	
	S. \ili{Nuaulu}	&	(w)ai, (w)a		&	\tsc{rel}	&	\tsc{AS}	&	\citet[61--63, 77--78, 164--170]{Bolton1990};\\	\hhline{~}
						&								&						&						&	\it{(w)ai} human; \it{(w)}a non-human	\\	
	\ili{Dobel}	&	ne	&	\tsc{rel}	&	\tsc{Aru}	&	\citet[171--177]{Hughes2000}	\\	
	\ili{Banda}	&	mana	&	‘which’	&	\ili{Banda}	&	\citet[564]{CollinsKaartinen1998}	\\	
	\ili{Kei} Kecil	&	en	&	\tsc{rel}	&	STB	&	\citet{Travis2019}	\\		\hhline{~}
		&	maŋ	&	role	&		&		\\	
	\ili{Kowiai}	&	ima	&	\tsc{rel}	&	STB	&	\citet[204, 230, 241]{WalkerWalker1991}	\\	
		\lspbottomrule
	\end{tabular}
\end{adjustbox}
\end{table}

\ssscp{Reduplication}{13-03-01-06}
Reduplication of many forms is widespread in the region involving
nouns, verbs, adverbs, numerals and quantifiers.
The specific functions of reduplication are tied to the part
of speech and verbal subclass in each language
(\citealp{Bolton1990,Grimes1991a,Grimes1991c,LaidigLaidig1991,
Hughes2000,Collins2003b,Coward2005,TaberTaber2015,Balle2017}
\citealp[116f]{Edwards2020}, \citealp{Tamelan2021,Verheijen1967}).
The forms and functions of reduplication in languages in the
region are consistent with forms
and functions of reduplication in languages west of the Blust
line, and cross linguistically around the world.
Most languages have multiple forms of reduplication,
as illustrated with data from \ili{Buru} in \trf{13-00}.

\begin{table}\tcp{Types of reduplication in Buru}{13-00}
\begin{adjustbox}{width=\textwidth}
	\st{2pt}
	\begin{tabular}{llllll}
\lsptoprule
Structure							&	\ili{Buru}								&	Gloss						&Function								&	Stem		&	Stem gloss	\\	\midrule
full stem \tsc{red}		&	emsian{\tl}em-sia-n	&	each 						&	[distributive]				&	emsian 	&	one (numeral)	\\
											&	emŋesa{\tl}em-ŋesa	&	together 				&	[plural]							&	emŋesa 	&	on an even plane	\\
\tsc{red} of root			&	geba{\tl}geba				&	various people  &	[variety]							&	geba 		&	person	\\
											&	gao{\tl}gao					&	grasp at 				&	[repeated]						&	gao 		&	hold	\\
											&	boho{\tl}boho				&	very 						&	[intensity,						&	boho 		&	bad, rotten, evil	\\
											&											&			 						&	lexicalised]					&			 		&					\\
											&	roho{\tl}roho				&	slow down				&	[degree, diminish]		&	roho 		&	reduce	\\
root-stem \tsc{red}		&	gosa{\tl}gosa-n			&	well-being 			&	[manner];							&	gosa 		&	good	\\
											&	sup{\tl}supa-k			&	do early in 		&	[degree];							&	supan 	&	morning	\\
											&											&	the morning 		&												&					&					\\
CVC-\tsc{red}					&	bal{\tl}bala				&	long ago 				&	[degree];							&	bala 		&	before	\\
CV-\tsc{red} various	&	gi{\tl}giwe					&	forcefully 			&	[intensity, manner];	&	giwe 		&	hard	\\
		\lspbottomrule
	\end{tabular}
\end{adjustbox}
\end{table}

No patterns of reduplication have yet been identified that show
potential as innovations for subgrouping purposes.
